\documentclass{article}
\usepackage{url}
\title{Lectures On Class Field Theory}
\author{Tamás Szamuely, notes by Francesco Minnocci}
\date{}

\usepackage{graphicx} % Required for inserting images
\usepackage[english]{babel}
\usepackage[a4paper]{geometry}
%\usepackage{fancyhdr}
\usepackage[position=top]{subfig}
%\usepackage{multirow}
%\usepackage{lipsum}
%\usepackage{titlesec}
%\usepackage{multicol}
\usepackage{setspace}
%\usepackage{enumitem}
%\usepackage{lipsum}
%\usepackage{rank-2-roots} %!!!!
\usepackage{calc}
\usepackage{enumerate}
\usepackage{centernot}
\usepackage{comment}
%\usepackage{microtype}
%\usepackage{todonotes}
%\usepackage{aeguill} %virgolette francesi
\usepackage{imakeidx}
\usepackage{faktor} %quozienti
\usepackage{stmaryrd} % simboli aggiuntivi


\usepackage{hyperref}
\hypersetup{
    colorlinks,
    citecolor=black,
    %filecolor=black,
    linkcolor=blue,
    urlcolor=black
}

% Teoremi e simboli
\usepackage{amsmath}
\usepackage{yhmath}
\usepackage{mathtools}
\usepackage{amssymb}
\usepackage{amsthm}
\usepackage{mathrsfs}
\usepackage{amsfonts}
%\usepackage{wasysym}
%\usepackage{thmtools}
% extensible arrows
\usepackage{extarrows}

%Box
\newsavebox{\boiteencadrement}
\newenvironment{encadrement}[1][gray!15]
{\par\addvspace{\topsep+10pt}%
    \def\couleurencadrement{#1}%
    \begin{lrbox}{\boiteencadrement}%
        \begin{minipage}{\linewidth-6\fboxsep-1pt}%
            }
            {\end{minipage}\end{lrbox}%
    \noindent\begin{tikzpicture}[
            inner sep=8pt,
            line width=0.75pt]
        \node[rounded corners = 5pt,
            draw,
            color=black,
            fill=\couleurencadrement] {\usebox{\boiteencadrement}};
    \end{tikzpicture}%
    \par\addvspace{\topsep+10pt}%
}

%Equazioni
\numberwithin{equation}{section}

%Teoremi
\theoremstyle{plain}

%\newtheorem{lem}[thm]{Lemma}
%\newtheorem{exer}[thm]{Esercizio}
%\newtheorem{cor}[thm]{Corollary}

\theoremstyle{definition}

\newcounter{thm}
\counterwithin{thm}{section}

%thm
\newtheorem{pretrm}[thm]{Theorem}
\newenvironment{thm}{\begin{encadrement}\begin{pretrm}}{\end{pretrm}\end{encadrement}}

%\newtheorem{defn}[thm]{Definition}
%\newcounter{defn}
%def
\newtheorem{predef}[thm]{Definition}
\newenvironment{defn}{\begin{encadrement}\begin{predef}}{\end{predef}\end{encadrement}}
%prop
\newtheorem{preprop}[thm]{Proposition}
\newenvironment{prop}{\begin{encadrement}\begin{preprop}}{\end{preprop}\end{encadrement}}
%cor
\newtheorem{precor}[thm]{Corollary}
\newenvironment{cor}{\begin{encadrement}\begin{precor}}
            {\end{precor}\end{encadrement}}
%lem
\newtheorem{prelem}[thm]{Lemma}
\newenvironment{lem}{\begin{encadrement}\begin{prelem}}{\end{prelem}\end{encadrement}}
%conj
\newtheorem{preconj}[thm]{Conjecture}
\newenvironment{conj}{\begin{encadrement}\begin{preconj}}{\end{preconj}\end{encadrement}}
%
\newtheorem{exmp}[thm]{Example}
\newtheorem{exer}[thm]{Exercise}
\newtheorem{rmk}[thm]{Remark}
\newtheorem*{rmk*}{Remark}
\newtheorem{facts}[thm]{Facts}
\newtheorem{fact}[thm]{Fact}
\newtheorem{constr}[thm]{Construction}
\newtheorem*{notat}{Notation}
\newtheorem*{note}{Note}


% Colonne
\usepackage{multicol}
\setlength\columnsep{20pt}
\setlength{\columnseprule}{0,5pt}

%% Comandi e operatori
\newcommand{\G}{\Gamma}

% No white line in equal arrows
\usepackage{tikz}
\usetikzlibrary{decorations.markings}
\tikzset{double line with arrow/.style args={#1,#2}{decorate,decoration={markings,%
                    mark=at position 0 with {\coordinate (ta-base-1) at (0,1pt);
                            \coordinate (ta-base-2) at (0,-1pt);},
                    mark=at position 1 with {\draw[#1] (ta-base-1) -- (0,1pt);
                            \draw[#2] (ta-base-2) -- (0,-1pt);
                        }}}}
\tikzset{Equal/.style={-,double line with arrow={-,-}}}

% mathbf
\newcommand{\N}{\mathbf{N}}
\newcommand{\Z}{\mathbf{Z}}
\newcommand{\Q}{\mathbf{Q}}
\newcommand{\QQ}{\overline{\mathbf{Q}}}
\newcommand{\R}{\mathbf{R}}
\newcommand{\K}{K}
\newcommand{\F}{\mathbf{F}}
\newcommand{\C}{\mathbf{C}}
\newcommand{\E}{\mathbf{E}}
% abelianization
\newcommand{\ab}{\text{ab}}


% mathcal
\newcommand{\A}{\mathbf{A}}
\renewcommand{\AA}{\mathcal{A}}
\newcommand{\Sn}{\mathcal{S}_n}
\newcommand{\D}{\mathcal{D}}

\newcommand{\OO}{\mathcal{O}}

% mathfrak
\newcommand{\g}{\mathfrak{g}}
\newcommand{\U}{\mathfrak{U}}
\newcommand{\gl}{\mathfrak{gl}}


\newcommand{\m}{\mathfrak{m}}

%For chapter 3
\newcommand{\cc}{\mathcal{c}}


\renewcommand{\sl}{\mathfrak{sl}}
\renewcommand{\sp}{\mathfrak{sp}}


% math operator
\DeclareMathOperator{\im}{im}
\DeclareMathOperator{\id}{id}
\DeclareMathOperator{\Hom}{Hom}
\DeclareMathOperator{\End}{End}
\DeclareMathOperator{\Aut}{Aut}
\DeclareMathOperator{\Inn}{Inn}
\DeclareMathOperator{\Out}{Out}
\DeclareMathOperator{\coker}{coker}
\DeclareMathOperator{\Stab}{Stab}
\DeclareMathOperator{\Gal}{Gal}
\DeclareMathOperator{\Frob}{Frob}
\DeclareMathOperator{\Span}{Span}
\DeclareMathOperator{\Sym}{Sym}
\DeclareMathOperator{\ad}{ad}
\DeclareMathOperator{\Sp}{Sp}
\DeclareMathOperator{\SO}{SO}
\DeclareMathOperator{\SL}{SL}
\DeclareMathOperator{\GL}{GL}
%\DeclareMathOperator{\OO}{O}
\DeclareMathOperator{\tr}{tr}
%new
\DeclareMathOperator{\car}{char}
\DeclareMathOperator{\Fib}{Fib}
%\DeclareMathOperator{\dep}{dep}
% \legendre{a}{p}
\newcommand{\legendre}[2]{\left(\frac{#1}{#2}\right)}
\renewcommand{\d}{\text{d}}

\newcommand\Iso{\xrightarrow{
        \,\smash{\raisebox{-0.65ex}{\ensuremath{\scriptstyle\sim}}}\,}}
\newcommand\restrict[1]{\raisebox{-.5ex}{$|$}_{#1}}

% Display math
\newcommand{\ssfrac}[2]{
    \raisebox{+0.3ex}{$#1$}
    /
    \raisebox{-0.3ex}{$#2$}
}
% Inline math
\newcommand{\sfrac}[2]{
    \raisebox{+0.3ex}{\scalebox{0.9}{$#1$}}
    /
    \raisebox{-0.3ex}{\scalebox{0.9}{$#2$}}
}

% boldsymbol

\renewcommand{\r}{{\boldsymbol{r}}}
\renewcommand{\k}{{k}}
\newcommand{\h}{{\boldsymbol{h}}}
\newcommand{\x}{{\boldsymbol{x}}}
\newcommand{\0}{{\boldsymbol{0}}}
\newcommand{\s}{{\boldsymbol{s}}}
\renewcommand{\t}{{\boldsymbol{t}}}
\newcommand{\e}{{\boldsymbol{e}}}
%map A_r
%\newcommand{\AR}{{\A}_{\r}}




\newcommand{\tn}{{\otimes n}}
\newcommand{\tens}{\otimes \ldots \otimes}
\newcommand{\n}[1]{\left\lVert #1 \right\lVert}
\newcommand{\res}[1]{\raisebox{-.5ex}{$|$}_{#1}}

%Elliptic curves

\renewcommand{\P}{\mathbf{P}}
\renewcommand{\S}{\operatorname{S}}
\newcommand{\OX}{\OO_{X,P}}
\DeclareMathOperator{\divv}{div}
\DeclareMathOperator{\Pic}{Pic}
\DeclareMathOperator{\Div}{Div}
\renewcommand{\L}{\mathcal{L}}
\newcommand\norm[1]{\left\lVert#1\right\rVert}
\newcommand{\p}{\mathfrak{p}}
\newcommand{\PP}{\mathfrak{P}}
\DeclareMathOperator{\Sel}{Sel}
\newcommand{\Selm}{\Sel^{(m)}}

%Galois groups and fundamental groups
\newcommand{\M}{\mathcal{M}}
% tr deg
\DeclareMathOperator{\trdeg}{tr.deg}
% spec
\DeclareMathOperator{\Spec}{Spec}
% frac
\DeclareMathOperator{\Frac}{Frac}


\begin{document}

\maketitle

\tableofcontents

\newpage

\section{Un peu d'histoire}
We start with some motivation: this course will be about number fields (i.e. finite extensions $K|\Q$) and their absolute Galois groups $G_K = \Gal(\overline{K}|K)$, where $\overline{K}$ is a fixed algebraic closure of $K$. Very little is known about the structure of such automorphism groups in general, and a large chunk of modern number theory/arithmetic geometry is devoted to studying them, e.g. via their representations. In some sense, class field theory consists in the study of its $1$-dimensional representations; more precisely, it describes the abelianization $G_K^{\ab} = G_K/[G_K,G_K]$ of the absolute Galois group (equivalently, of its abelian quotients).

\subsection{Reciprocity laws}
Let $p$ be a prime number, and let $a$ be an integer coprime to $p$. The \emph{Legendre symbol} $\legendre{a}{p}$ is defined as
\[
	\legendre{a}{p} = \begin{cases}
		1  & \text{if $a$ is a square modulo $p$},     \\
		-1 & \text{if $a$ is not a square modulo $p$}.
	\end{cases}
\]
\textbf{Euler} proved that $\legendre{a}{p} = a^{\frac{p-1}{2}} \bmod p$, and \textbf{Gauss} proved the following law:
\begin{thm}[Quadratic reciprocity]
	Let $a\ne b$ be different odd primes. Then,
	\[
		\legendre {a}{b}\cdot\legendre{b}{a} = (-1)^{\frac{a-1}{2}\frac{b-1}{2}}.
	\]
	Morever, the following subsidiary laws hold:
	\begin{enumerate}
		\item $\legendre{-1}{p} = (-1)^{\frac{p-1}{2}}$;
		\item $\legendre{2}{p} = (-1)^{\frac{p^2-1}{8}}$.
	\end{enumerate}
\end{thm}
These laws, together with its multiplicativity in the parameter $a$, make it possible to compute the Legendre symbol.

Another way of thinking about class field theory is as an attempt to define a ``very general'' reciprocity law, from which the quadratic reciprocity law and its higher degree analogues can be deduced. One such generalization is \textbf{Hilbert}'s reciprocity law: recall that $\Q$ embeds into the $p$-adic field $\Q_p$ for every prime $p$, and also into the real numbers $\Q_\infty\coloneqq\R$.
For every prime $p$ (including $p=\infty$) and for all $a,b\in \Q^\times$, we define the \emph{Hilbert symbol} $(a,b)_p$ as
\[
	(a,b)_p = \begin{cases}
		1  & \text{if $ax^2 + by^2 = 1$ has a solution in $\Q_p^2$}, \\
		-1 & \text{otherwise}.
	\end{cases} asd
\]

Recall that the norm $N:\Q_p(\sqrt{a})\to \Q_p$ is defined as $N(x+y\sqrt{a}) = x^2 - ay^2$ if $a\notin \Q_p^{\times 2}$, and as $N=\id$ otherwise. More generally, if $L|K$ is a finite Galois extension with group $G$, the norm $N_{L|K}:L\to K$ is
\[
	N(x)=\prod_{\sigma\in G} \sigma(x)
\]
\begin{lem}
	The following statements are equivalent:
	\begin{enumerate}
		\item $(a,b)_p = 1$;
		\item The quadratic equation $z^2 = ax^2 + by^2$ has a non-trivial solution;
		\item $a\in \Q_p$ is the norm of an element of $\Q_p(\sqrt{b})$.
		\item $b\in \Q_p$ is the norm of an element of $\Q_p(\sqrt{a})$.
	\end{enumerate}
\end{lem}
\begin{proof}
	The equivalence $1\iff 2$ is by definition; and $3\iff 4$ is clear.
	For $2\iff 3$, if $b\in \Q_p^{\times 2}$, then we are done. Otherwise, if $b\notin \Q_p^{\times 2}$, then $N(\alpha+\beta\sqrt{b})= \alpha^2 - b\beta^2$, and if $a\coloneqq \alpha^2 - b\beta^2$ we see that $z^2 = ax^2 + by^2$ has a non-trivial solution, and vice versa.
\end{proof}
For $p\ne 2, \infty$ and $a,b\in\Z$, write
$a=p^{v_p(a)}u$ and $b=p^{v_p(b)}v$ with $p\nmid uv$. The relation between the Hilbert symbol and the Legendre symbol is
\begin{align}
	(a,b)_p
	 & = (-1)^{\frac{p-1}{2}v_p(a)v_p(b)}
	\legendre{u}{p}^{v_p(b)}\legendre{v}{p}^{v_p(a)}                    \\
	 & = \overline{(-1)^{v_p(a)v_p(b)}a^{v_p(b)}b^{-v_p(a)}}^{(p-1)/2}.
\end{align}
where the bar denotes reduction modulo $p$.
In particular, if $p$ does not divide $a$, then $(a,b)_p = \legendre{a}{p}^{v_p(b)}$ coincides with the Legendre symbol.
\begin{thm}[Hilbert reciprocity law]
	For all $a,b\in \Q^\times$, we have
	\[
		\prod_{p} (a,b)_p = 1
	\]
\end{thm}
As for computation, one shows that if $p$ does not divide $2ab$ (including $p=\infty$ for $a,b>0$), then $(a,b)_p = 1$. Moreover, $(a,b)_2=(-1)^{\frac{a-1}{2}\frac{b-1}{2}}$ for odd $a,b$. From this and multiplicativity in the first variable, we immediately deduce the quadratic reciprocity law by taking $a,b$ to be odd primes.

An advantage of this formulation is that it generalizes to number fields $K$, by considering the completions of $K$ at all of its (finite and infinite) ``places''. Hilbert then went on to pose his 9th problem, in which he asked for a generalization of this law to $n$-th powers in arbitrary number fields; this is what class field theory achieves. We first consider the case of extensions which contain a primitive $n$-th root of unity:

\subsection{Kummer case}
Assume now $K|\Q$ is a finite extension such that $\mu_n\subset K$, and let $\OO_K$ be its ring of integers. If $P\subset \OO_K$ is a prime ideal, we would like to define a symbol $(a,b)_P\in \mu_n$ such that a product formula holds, i.e.
\[
	\prod_{P} (a,b)_P = 1
\]
if $a,b\in K^\times$. Moreover, for $P$ finite (which is always the case for $n\ne 2$) which does not divide $(n)$ we shall have \[
	(a,b)_P = \overline{(-1)^{v_P(a)v_P(b)}a^{v_P(b)}b^{-v_P(a)}}^{(q-1)/n},\] where $q = \#(\OO_K/P)$. In particular, if $P$ does not divide $(a)$, then $(a,b)_P =
	a^{\frac{q-1}{n}v_P(b)}\bmod P$.

So, if $v_P(b)=1$, then
\[
	\legendre{a}{P} \coloneqq (a,b)_P
\]
does not depend on the choice of $b$ and it is the unique $n$-th root of unity, that is $a^{\frac{q-1}{n}}\bmod P$; in other words, $\legendre{a}{P} = 1$ if and only if $x^n =a \bmod P$ has a solution. Some more properties of this symbol are:
\begin{enumerate}
	\item $(a,b)_P=1 \iff$ $b$ is a norm from $K_P(\sqrt[n]{a}) \iff$ $a$ is a norm from $K_P(\sqrt[n]{b})$;
	\item $(a,b)_P$ is multiplicative in both variables and anticommutative;
\end{enumerate}

When $\mu_n$ is \textit{not} contained in $K_P$, an ``Hilbert'' symbol still exists but is longer associated to $a,b\in K^\times$. Suppose for a moment that $\mu_n\subset K_P$. Then, the extension $K_P(\sqrt[n]{a})|K_P$ is Galois for all $a\notin K_P^{\times n}$, and its Galois group is
\[
	G_a\coloneqq\Gal(K_P(\sqrt[n]{a})|K_P) \cong \Z/n\Z
\]
Then, the group $ \Hom(G_a,\mu_n)\cong \Z/n\Z$ is generated by the character
\[
	\overline{\chi_a}:\sigma\mapsto \frac{\sigma(\sqrt[n]{a})}{\sqrt[n]{a}}\in \mu_n \cong \Z/n\Z.
\]
In general, we shall construct a generalized Hilbert symbol
\begin{align*}
	\Hom\left(\Gal(\overline{K_P}|K_P),\Z/n\Z\right) \times K_P^\times & \to \Z/n\Z,        \\
	[\chi,b]                                                           & \mapsto (\chi,b)_P
\end{align*}

and we'll thus be able to prove the \textbf{Artin} reciprocity law: fix some $b\in K^\times$, and suppose that $\chi$ is a character in $\Hom(G_K,\Z/n\Z)$. Then
\[
	\sum_P (\chi_P,b)=0,
\]
where $\chi_P$ is the restriction of $\chi$ to $G_{K_P}^\ab$ (recall that it embeds into $G_K^\ab$).

The generalized Hilbert symbol defines a reciprocity map
\[
	K_P^\times \to G_{K_P}^\ab,
\]
which is ``almost'' an isomorphism. In particular, for every $n$ we'll have an isomorphism of finite groups
\[
	K_P^\times/n \to G_{K_P}^\ab/n.
\]
This is one of the remarkable features of class field theory, giving a description of the abelian extensions of a field in terms of its arithmetic.
To define such a map, it's enough to define the dual map
\begin{align*}
	\Hom(G_{K_P}^\ab/n,\Z/n\Z) & \to \Hom(K_P^\times/n,\Z/n\Z), \\
	[\chi]                     & \mapsto [b\mapsto (\chi,b)_P].
\end{align*}

\section{Computing the Hilbert symbol}

Recall the definition: for $a,b\in\Q_p^\times$,
\[
	(a,b)_p = \begin{cases}
		1  & \text{if $ax^2+by^2=1$ is solvable in $\Q_p$}, \\
		-1 & \text{otherwise}.
	\end{cases}
\]
We have seen that
\[
	(a,b)_p = 1 \iff z^2 - ax^2 - by^2 \text{ has a non-trivial solution in } \Q_p^3 \iff b \text{ is a norm from } \Q_p(\sqrt{a}),
\]
where $N(x+y\sqrt{a}) = x^2 - ay^2$.

\begin{thm}[Hilbert reciprocity law]\label{thm:hilbertrecip}
	For all $a,b\in\Q^\times$,
	\[
		\prod_{p} (a,b)_p = 1,
	\]
	where the product runs over all primes $p$, including $p=\infty$.
\end{thm}
Note that the product is finite: $(a,b)_p=1$ for all $p\nmid 2ab$ (see Lemma~\ref{lem:units} below).

\begin{lem}\label{lem:basic}
	Let $a,b,c\in\Q_p^\times$. Then
	\begin{enumerate}
		\item $(1,c)_p = 1$;
		\item $(-c,c)_p = 1$;
		\item if $(a,c)_p = 1$, then $(ab,c)_p = (a,c)_p(b,c)_p$ (this holds in general, see the Corollary below);
		\item $(c,c)_p = (-1,c)_p$.
	\end{enumerate}
\end{lem}
\begin{proof}
	\begin{enumerate}
		\item Consider $N:\Q_p(\sqrt{c})\to\Q_p$. Then $N(1)=1$.
		\item If $c\in\Q_p^{\times 2}$, then $\Q_p(\sqrt{c})=\Q_p$, the norm is the identity and $N(-c)=-c$. If $c\notin\Q_p^{\times 2}$, then $N(\sqrt{c}) = -c$. In both cases $-c$ is a norm from $\Q_p(\sqrt{c})$.
		\item If $a=N(x)$ and $b=N(y)$, then $ab=N(x)N(y)=N(xy)$, since the norm is multiplicative by definition.

		      Moreover, if $a=N(x)$, then $a^{-1}=N(x^{-1})$; if in addition $ab=N(z)$, then
		      \[
			      b=a^{-1}(ab)=N(x^{-1})N(z)=N(x^{-1}z).
		      \]
		      So, if $(a,c)_p=1$ (i.e.\ $a=N(x)$ is a norm from $\Q_p(\sqrt{c})$), then $(b,c)_p=1\iff(ab,c)_p=1$.
		\item By (2) and (3),
		      \[
			      (c,c)_p = (-c,c)_p(c,c)_p = (-c^2,c)_p = (-1,c)_p,
		      \]
		      because in general $(ad^2,b)_p = (a,b)_p$.\qedhere
	\end{enumerate}
\end{proof}
\begin{rmk}
	Since $(ad^2,b)_p=(a,b)_p$, the Hilbert symbol is a well-defined map
	\[
		(\,\cdot\,,\cdot\,)_p : \left(\Q_p^\times/\Q_p^{\times 2}\right)\times\left(\Q_p^\times/\Q_p^{\times 2}\right)\to\mu_2.
	\]
\end{rmk}

\begin{lem}\label{lem:units}
	Let $p>2$. Then $(u,v)_p=1$ for all $u,v\in\Z_p^\times$.
	(Recall that $\Z_p$ is local and $\Z_p^\times=\Z_p\setminus(p)$.)
\end{lem}
\begin{proof}
	Let $\overline{u}\coloneqq u\bmod p$ and $\overline{v}\coloneqq v\bmod p$; they are non-zero elements of $\F_p$ (units of the local ring).
	We first show that $1-\overline{u}x^2-\overline{v}y^2$ has a solution in $\F_p$. Since $\F_p^\times$ is cyclic of order $p-1$,
	\[
		\#\{x^2 : x\in\F_p\} = \frac{p+1}{2}.
	\]
	Similarly, if $S_{\overline{u}}\coloneqq\{\overline{u}x^2:x\in\F_p\}$ and $S_{\overline{v}}\coloneqq\{1-\overline{v}y^2 : y\in\F_p\}$, then $\#S_{\overline{u}}=\#S_{\overline{v}}=
		\frac{p+1}{2}$.
	Thus, $S_{\overline u}\cap S_{\overline v}\neq\emptyset$ by the

	We conclude by Hensel's lemma:
\end{proof}

\begin{lem}[Hensel]
	Let $f\in\Z_p[X]$ and $a_0\in\Z_p$ with $f(a_0)\equiv 0\mod p$ and $f'(a_0)\not\equiv 0\mod p$. Then there exists $a\in\Z_p$ with $f(a)=0$ and $a\equiv a_0\mod p$.
\end{lem}
This allows us to lift a solution mod $p$ to a solution in $\Z_p$ (to apply it in the proof above, we have to fix one variable).

\begin{prop}\label{prop:oddp}
	Let $p>2$ and $a,b\in\Q_p^\times$, with $a=p^\alpha u$, $b=p^\beta v$ and $u,v\in\Z_p^\times$. Then
	\[
		(a,b)_p = (-1)^{\alpha\beta\frac{p-1}{2}}\legendre{\overline{u}}{p}^{-\beta}\legendre{\overline{v}}{p}^{\alpha}.
	\]
\end{prop}
\begin{proof}
	Since $(a,b)_p$ depends only on $a,b \bmod \Q_p^{\times 2}$, we may assume $\alpha,\beta\in\{0,1\}$.
	\begin{itemize}
		\item $\alpha=\beta=0$: this is Lemma~\ref{lem:units}.
		\item $\alpha=1$, $\beta=0$: we need $(pu,v)_p=\legendre{\overline{v}}{p}$. By Lemma~\ref{lem:units} we have $(u^{-1},v)_p=1$, hence by Lemma~\ref{lem:basic}(3)
		      \[
			      (pu,v)_p = (pu,v)_p\underbrace{(u^{-1},v)_p}_{=1} = (p,v)_p.
		      \]
		      If $v\in\Q_p^{\times 2}$, then $(p,v)_p=(p,1)_p=1=\legendre{\overline{v}}{p}$ (note that $\overline{v}$ is then also a square mod $p$).
		      If $v\notin\Q_p^{\times 2}$, then $(p,v)_p=-1$, because $z^2-px^2-vy^2$ has no non-trivial solution mod $p$, hence none in $\Q_p$ by homogeneity; also $\legendre{\overline{v}}{p}=-1$ (by Hensel, since $v\notin\Q_p^{\times 2}$).
		\item $\alpha=\beta=1$: we need $(pu,pv)_p=(-1)^{\frac{p-1}{2}}\legendre{\overline{u}}{p}\legendre{\overline{v}}{p}$. By Lemma~\ref{lem:basic}(2) we have $(-pv,pv)_p=1$, so by Lemma~\ref{lem:basic}(3)
		      \[
			      (pu,pv)_p = (pu,pv)_p\underbrace{(-pv,pv)_p}_{=1} = (-p^2uv,pv)_p = (-uv,pv)_p = (pv,-uv)_p,
		      \]
		      and now we are in the previous case, which gives
		      \[
			      (pv,-uv)_p = \legendre{\overline{-uv}}{p} = \legendre{-1}{p}\legendre{\overline{u}}{p}\legendre{\overline{v}}{p},
		      \]
		      where $\legendre{-1}{p}=(-1)^{\frac{p-1}{2}}$, and we used that the Legendre symbol is multiplicative.\qedhere
	\end{itemize}
\end{proof}

\begin{cor}
	Let $p>2$. Then $(a,c)_p(b,c)_p=(ab,c)_p$ for all $a,b,c\in\Q_p^\times$. In particular, $(\,\cdot\,,\cdot\,)_p$ is bi-multiplicative.
\end{cor}
\begin{proof}
	Write $a=p^\alpha u$, $b=p^\beta v$ and $c=p^\gamma w$ with $u,v,w\in\Z_p^\times$. By Proposition~\ref{prop:oddp},
	\begin{align*}
		(a,c)_p(b,c)_p
		 & = (-1)^{\frac{p-1}{2}\alpha\gamma}\legendre{\overline{u}}{p}^{\gamma}\legendre{\overline{w}}{p}^{\alpha}\cdot(-1)^{\frac{p-1}{2}\beta\gamma}\legendre{\overline{v}}{p}^{\gamma}\legendre{\overline{w}}{p}^{\beta} \\
		 & = (-1)^{\frac{p-1}{2}\gamma(\alpha+\beta)}\legendre{\overline{uv}}{p}^{\gamma}\legendre{\overline{w}}{p}^{\alpha+\beta} = (ab,c)_p.\qedhere
	\end{align*}
\end{proof}

\paragraph{Case $p=2$.}
If $a=2^\alpha u$ and $b=2^\beta v$ with $u,v\in\Z_2^\times$, define
\[
	\varepsilon(u)\coloneqq\frac{u-1}{2}\bmod 2,\qquad \omega(u)\coloneqq\frac{u^2-1}{8}\bmod 2.
\]
Then
\[
	(a,b)_2 = (-1)^{\varepsilon(u)\varepsilon(v)+\alpha\omega(v)+\beta\omega(u)}.
\]
(Proof not given.) In particular, $(\,\cdot\,,\cdot\,)_p$ is bi-multiplicative for every $p$.
\begin{proof}[Proof of Theorem \ref{thm:hilbertrecip}]
	Using bi-multiplicativity, we reduce to the case where each of $a,b$ is either $-1$ or a prime.
	\begin{itemize}
		\item Case $a=-1$: homework.
		\item Case $a=b=q$ prime: $(q,q)_p=(-1,q)_p$ by Lemma~\ref{lem:basic}(4), so we are reduced to the previous case.
		\item Case $a=2$, $b>2$ prime: $(2,b)_p=1$ for $p\notin\{2,b\}$ (also for $p=\infty$, as $2,b>0$). Moreover, by the $p=2$ formula,
		      \[
			      (2,b)_2=(-1)^{\omega(b)},
		      \]
		      and by Proposition~\ref{prop:oddp}
		      \[
			      (2,b)_b=\legendre{2}{b}=(-1)^{\frac{b^2-1}{8}}.
		      \]
		      The product of these two factors is $1$.
		\item Case $a\neq b$ odd primes: $(a,b)_p=1$ for $p\notin\{2,a,b\}$ (also for $p=\infty$), and
		      \[
			      (a,b)_p = \begin{cases}
				      (-1)^{\frac{a-1}{2}\frac{b-1}{2}} & p=2, \\[2pt]
				      \legendre{a}{b}                   & p=b, \\[2pt]
				      \legendre{b}{a}                   & p=a.
			      \end{cases}
		      \]
		      Hence the product is $1$ by quadratic reciprocity.\qedhere
	\end{itemize}
\end{proof}

\section{Cohomology of groups}
Already for the local part of theory we'll require some input from homological algebra, mainly the description of the first cohomology group $H^1(G,A)$ for $G$ some Galois group, and its interpretation in terms of crossed homomorphisms. A concise reference for this material can be found in the notes \cite{homalg}. We recall some results here for the reader's convenience:

\begin{thm}[Cohomology of cyclic groups]
	Let $G$ be a finite cyclic group of order $n$, generated by $\sigma$. Then,
	\[
		\begin{cases}
			H^0(G,A) = A^G,           \\[2pt]
			H^1(G,A) = A/(1-\sigma)A, \\[2pt]
			H^2(G,A) = A^G/NA,        \\[2pt]
			H^i(G,A) \cong H^{i+2}(G,A) \text{ for all } i\ge 1,
		\end{cases}
	\]
\end{thm}
This comes out of the ``periodic'' resolution of $\Z$ as a trivial $G$-module:
\[
	\ldots \xrightarrow{\sigma-1} \Z[G] \xrightarrow{N} \Z[G] \xrightarrow{\sigma-1} \Z[G] \xrightarrow{N} \Z[G] \xrightarrow{\varepsilon} \Z \to 0,
\]
where $N = 1+\sigma+\ldots+\sigma^{n-1}$ is the norm map and $\varepsilon$ sends $\sigma$ to $1$.
A similar computation holds for the infinite cyclic group $G=\Z$:
\[
	H^0(\Z,A) = A^\Z,\qquad H^1(\Z,A) = A/(1-\sigma)A,\qquad H^i(\Z,A) = 0 \text{ for all } i\ge 2.
\]
For a subgroup $H\subset G$, we denote by
\[M^G_H(A)\coloneqq  \Hom_H(\Z[G],A)\] the $G$-module with action
\[
	\sigma(\phi)(g) = \phi(\sigma^{-1}g),\qquad \sigma\in G, g\in G, \phi\in M^G_H(A).
\]
Then, \textit{Shapiro}'s lemma tells us that there is a canonical isomorphism
\begin{equation}\label{eq:shapiro}
	H^i(G,M^G_H(A)) \cong H^i(H,A).
\end{equation}


\begin{thebibliography}{9}
	\bibitem{milne} James S. Milne, \emph{Class field theory}, \url{https://www.jmilne.org/math/CourseNotes/CFT.pdf}.
	\bibitem{homalg} Tam\'as Szamuely, \emph{Notes on homological algebra}, \url{https://pagine.dm.unipi.it/tamas/ist-alg2new.pdf}.
\end{thebibliography}

\end{document}
